
% \newpage
% \section*{Limitation}
% Limitation

% \section*{Ethics Consideration}

% In this paper, we investigate the vulnerabilities of \rag, a retrieval-augmented generation paradigm, to adversarial poisoning attacks. Our work aims to uncover both the strengths and weaknesses of \rag's security mechanisms, with the goal of improving its robustness and reliability in real-world applications. While our research contributes to the advancement of secure AI systems, it also raises important ethical considerations.

% \textbf{Stakeholder Considerations.}  
% The primary stakeholders affected by this research include developers and users of \rag systems, as well as the broader community relying on these systems for accurate and reliable information retrieval. By exposing vulnerabilities in \rag, our work may initially undermine trust in its security. However, this transparency is essential for driving improvements in the system's defenses, ultimately leading to more secure and trustworthy AI applications. Additionally, our findings could inform policymakers and practitioners about the risks associated with deploying \rag in critical domains such as healthcare, cybersecurity, and geographic information systems.

% \textbf{Potential Harms and Benefits.}  
% The primary benefit of this research is the identification of vulnerabilities in \rag, which enables a significant warning of this new tool and encourages the development of more robust defenses against poisoning attacks. However, there is a potential harm in publicly disclosing these vulnerabilities, as malicious actors could use this information to craft more sophisticated attacks. To balance this, we try to discuss some potential countermeasures for improving \rag's security even if it's not enough.

% \textbf{Future Research and Mitigation.}  
% Our study highlights the need for tailored defenses against poisoning attacks in \rag systems. We propose several potential countermeasures, such as filtering poisoning content and validating the consistency of answers by CoTs. Future research should explore these directions to ensure the safe deployment of \rag in real-world applications.

% \section*{Open Science Policy}

% In alignment with the principles of open science, we are committed to transparency and reproducibility in our research. To this end, we will make the following artifacts publicly available:

% \begin{itemize}
%     \item \textbf{Datasets:} All datasets used in our experiments, including geographic, medical, and cybersecurity datasets, will be provided in anonymized form to ensure privacy and compliance with ethical guidelines.
%     \item \textbf{Code and Scripts:} We will release the source code for our experiments, including the implementation of \attack and the evaluation framework for \rag systems. This enables other researchers to replicate our findings and build upon our work.
%     \item \textbf{Documentation:} Detailed documentation will be provided to guide users in setting up and running the experiments, ensuring accessibility for the broader research community.
% \end{itemize}

% All research artifacts are available at: \url{https://anonymous.4open.science/r/\rag-Poisoning}.

% By sharing these resources, we aim to foster collaboration and accelerate progress in the field of secure retrieval-augmented generation systems. We encourage researchers to use these materials responsibly and to contribute to the development of robust defenses against adversarial attacks.



% \section*{Ethics Consideration}
% %\vspace{-0.1cm}

% In this paper, we conduct a systematic study of state-of-the-art LLM watermarkers, focusing on their robustness against watermark removal attacks. We aim to understand how various design choices affect watermarkers' resilience and identify best practices for operating watermarkers in adversarial environments. 


% {\bf Stakeholder Considerations.} The primary stakeholders affected by this research include users and developers of LLM watermarkers, as well as the broader community relying on these techniques. By identifying the strengths/limitations of various watermarkers, our work could influence the perceived reliability and deployment of LLM watermarkers in critical applications. Conversely, exposing these vulnerabilities allows for the improvement of the current techniques, ultimately contributing to more secure and robust systems.

% {\bf Potential Harms and Benefits.}
% Exposing vulnerabilities in widely adopted security mechanisms can yield contrasting outcomes. Initially, it may erode trust in LLM watermarkers as reliable safeguards, potentially leading to their underuse and leaving systems exposed to unchecked risks. However, by illuminating these weaknesses, our research catalyzes the development of more robust techniques and fosters a nuanced comprehension of their constraints, ultimately fortifying the long-term security ecosystem.

% {\bf Future Research and Mitigation.}
% Our study has identified several potential countermeasures to address the vulnerabilities we uncovered. These recommendations are designed to steer future research and assist developers in bolstering the security of LLM watermarkers.

% \section*{Open Science Policy}

% In compliance with USENIX Security's new open science policy, we have released all relevant research artifacts, including datasets, scripts, and source code: \url{https://anonymous.4open.science/r/WaterPark}.
 

% We evaluate eleven watermarking schemes
% from related work and empirically determine their decision
% thresholds for the CIFAR-10 and ImageNet datasets. Then, we
% measured the performance of a large set of removal attacks
% against all watermarking schemes and ablate over multiple
% parameters for each scheme and removal attack. We use the
% Nash equilibrium to evaluate a scheme’s robustness against (i)
% all attacks, (ii) categories of attacks, and (iii) single attacks.
% Our results show that none of the schemes is robust against
% all attacks. We break down these results by analyzing each
% attack category’s effectiveness and find that the most effective
% removal attack category are model extraction attacks, followed
% by model modification attacks. We show that transfer learning
% removes all watermarks on CIFAR-10, but there exists no such
% dominant attack for ImageNet. We create a combined attack
% composed of (1) transfer learning and (2) label smoothing
% that removes all eleven watermarks. Finally, we discuss the
% practicality of the removal attacks, e.g., their monetary costs
% and the dataset availability of the attacker and propose guidelines for evaluating the robustness of DNN watermarking. We
% hope that our work will improve future evaluations of DNN
% watermarking schemes.